\chapter{Conclusions and future work}

\section{Conclusions}
The development of Micorrizas has made it possible to create a functional Android serious game based on a real educational activity. The project has achieved its main technical objective by producing a mobile application that combines geolocation, cooperative multiplayer interaction, mission-based gameplay and a final group result. These elements work together to support an outdoor learning experience in the Jardín de los Sentidos.

One of the main conclusions of the project is that the game is not simply a gamified version of an existing activity. During development, it became clear that the original visit was being used as the basis for designing a serious game with its own structure, narrative, mechanics and progression. The educational activity provided the context, contents and objectives, but the final result required specific design decisions in order to transform that material into an interactive experience.

The project has also shown the value of connecting mobile gameplay with a real environment. By using the garden as the playable space, the game encourages players to observe, move, communicate and make decisions in relation to the physical surroundings. This supports the objective of promoting awareness and appreciation of biodiversity through interaction with the natural environment.

However, the educational impact of this objective cannot yet be fully measured. The application is prepared to be introduced as a tool in the subject, and the lecturers have expressed the intention of using it in a real teaching session during the next academic year. Once the game is used with students in that context, it will be possible to evaluate more precisely whether the experience increases their awareness, appreciation and understanding of the biodiversity of the Jardín de los Sentidos.

Working on mycorrhizae has enabled me to put much of the knowledge I acquired during my degree into practice, as well as to incorporate new knowledge that has proved necessary during the course of the project.

Overall, the result achieved is satisfactory, as the game meets the project’s expectations and demonstrates the feasibility of building a geolocation-based cooperative experience in Unity.


\section{Future work}
The main line of future work is the use of the application in a real session of the subject. This would allow the game to be tested with its intended users and under authentic teaching conditions. The results of that implementation could be used to evaluate the educational impact of the game, especially in relation to student engagement and collaborative learning.

A future evaluation could combine several methods, such as observation during the activity, questionnaires before and after the session, teacher feedback and analysis of the final group scores. This would make it possible to assess not only whether the application works technically, but also whether it contributes to the learning objectives of the original activity.

From a technical perspective, future work could focus on improving GPS accuracy and robustness during outdoor use. Although the prototype already includes geolocation-based mission activation, further testing with larger groups and different mobile devices would help to refine activation areas, GPS indicators and tolerance margins.

The content of the game could also be expanded. New plant species, sounds, images and questions could be added to enrich the missions and make the experience more varied. Additional missions could also be designed for other botanical or educational spaces, allowing the structure of Micorrizas to be reused beyond the Jardín de los Sentidos.

Finally, the teacher-facing part of the application could be improved. Future versions could include tools to export results, review group performance or configure the content of the missions before the session. These improvements would make the application more useful as a teaching resource and would support its integration into the subject in a more stable way.
